\documentclass{article}

\usepackage[margin=1in]{geometry}

\usepackage{parskip}

\begin{document}

% Page 1

Nature of S‑States in the Oxygen-Evolving Complex Resolved by
High-Energy Resolution Fluorescence Detected X‑ray Absorption
Spectroscopy
Maria Chrysina, Maria Drosou, Rebeca G. Castillo, Michael Reus, Frank Neese, Vera Krewald,
Dimitrios A. Pantazis,* and Serena DeBeer*
Cite This: J. Am. Chem. Soc. 2023, 145, 25579−25594
Read Online
ACCESS
Metrics \& More
Article Recommendations
*
sı
Supporting Information
ABSTRACT: Photosystem II, the water splitting enzyme of photosynthesis,
utilizes the energy of sunlight to drive the four-electron oxidation of water to
dioxygen at the oxygen-evolving complex (OEC). The OEC harbors a
Mn4CaO5 cluster that cycles through five oxidation states Si (i = 0−4). The
S3 state is the last metastable state before the O2 evolution. Its electronic
structure and nature of the S2 → S3 transition are key topics of persisting
controversy. Most spectroscopic studies suggest that the S3 state consists of
four Mn(IV) ions, compared to the Mn(III)Mn(IV)3 of the S2 state.
However, recent crystallographic data have received conflicting interpreta-
tions, suggesting either metal- or ligand-based oxidation, the latter leading to
an oxyl radical or a peroxo moiety in the S3 state. Herein, we utilize high-
energy resolution fluorescence detected (HERFD) X-ray absorption spectroscopy to obtain a highly resolved description of the Mn
K pre-edge region for all S-states, paying special attention to use chemically unperturbed S3 state samples. In combination with
quantum chemical calculations, we achieve assignment of specific spectroscopic features to geometric and electronic structures for all
S-states. These data are used to confidently discriminate between the various suggestions concerning the electronic structure and the
nature of oxidation events in all observable catalytic intermediates of the OEC. Our results do not support the presence of either
peroxo or oxyl in the active configuration of the S3 state. This establishes Mn-centered storage of oxidative equivalents in all
observable catalytic transitions and constrains the onset of the O−O bond formation until after the final light-driven oxidation event.
■ INTRODUCTION
Photosystem II (PSII) of plants, algae, and cyanobacteria
converts sunlight into chemical energy in the form of “high
energy” electrons derived from water. These electrons are used
in CO2 reduction to glucose, which is the primary form of food
for all living organisms. PSII is thus a key enzyme in the flow of
energy into the biosphere, and its unique ability to oxidize
water makes its study of paramount importance both for basic
research and as the only model system for developing artificial
water-splitting catalysts. Water splitting takes place at the
Mn4CaO5 cluster of the oxygen-evolving complex (OEC)
(Figure 1a), embedded in the protein scaffold of PSII.1−5
The cluster cycles through a progression of four light-
induced sequential oxidations, S0 → S1 → S2 → S3 → [S4],
mediated by redox active tyrosine residue Tyr161 (YZ).
Dioxygen is formed and released during the last S3 → [S4] →
S0 transition, where S4 is a transient state (Figure 1b). The
catalytic cycle includes two water binding events, probably
during or after the S2 → S3 transition and prior to
reconstitution of the S0 state, but the details of water binding
as well as the position and protonation states of the substrates
for dioxygen formation remain uncertain. Commonly accepted
formal oxidation states of the four Mn ions in the S0, S1, and S2
states are Mn(III)3Mn(IV), Mn(III)2Mn(IV)2, and Mn(III)-
Mn(IV)3, respectively.6−11 Therefore, the S0−S2 oxidation
events are unanimously considered to be Mn-based. However,
the localization of the S2 → S3 oxidation remains a
fundamental mechanistic question since differing suggestions
exist about the nature of the transition as well as the electronic
structure of the S3-state itself.10,12−18
All interpretations of electron paramagnetic resonance
(EPR) observations for all different signals associated with
the S3 state are consistent with Mn-centered oxidation in the S2
→ S3 transition, leading to a Mn(IV)4 assignment for the S3
state.19−25 However, in X-ray absorption near edge spectros-
copy (XANES) experiments, the shift of the Mn K-edge to
higher energy in the S2 → S3 transition was found to be smaller
than that in the S0 → S1 and S1 → S2 transitions. This
Received:
June 9, 2023
Revised:
October 13, 2023
Accepted:
October 13, 2023
Published: November 16, 2023
Article
pubs.acs.org/JACS
© 2023 The Authors. Published by
American Chemical Society
25579
https://doi.org/10.1021/jacs.3c06046
J. Am. Chem. Soc. 2023, 145, 25579−25594
This article is licensed under CC-BY 4.0
Downloaded via FREIE UNIV BERLIN on November 28, 2025 at 14:20:54 (UTC).
See https://pubs.acs.org/sharingguidelines for options on how to legitimately share published articles.

\newpage

% Page 2

observation along with Kβ X-ray emission spectra during the S2
→ S3 transition9,15,26 has received two alternative interpreta-
tions: either ligand-based oxidation during this transition,15 or
a coordination sphere change from a five-coordinate Mn(III)
ion to an octahedral Mn(IV) ion10,17 as a result of water
binding (Figure 2, oxo-hydroxo). The first scenario would
mean that the Mn oxidation states in the S3 state remain
Mn(III)Mn(IV)3 and that a ligand radical species is formed
instead (Figure 2, oxyl-oxo). This scenario appears consistent
with two-flash (2F) structural models derived by recent serial
femtosecond X-ray free electron laser crystallography (SFX-
XFEL),27−30 Crystallographic data have received conflicting
interpretations, ranging from an all-Mn(IV) oxo-hydroxo,29−31
to a ligand oxidized oxyl-oxo or peroxo species (Figure
2).27,28,32−34
The above distinct formulations for the S3 state have
profound implications for the water oxidation mechanism. The
possibility of oxygen-radical formation after the second flash
(2F) led to the hypothesis of an “early-onset” O−O bond
formation already from the S3 state.18,32,36−40 By contrast, if
the S3 state is Mn(IV)4, the O−O bond formation must take
place after the last flash of the cycle (3F), either in a transient
S4 state or after formation of the S3YZ
• intermediate.
Computationally, all three types of models have been proposed
as part of possible water oxidation mechanisms.32,39,41,42
However, the relative energies between computational models
of the different S3 variants and the correspondence to the
spectroscopic observations remain inconclusive.
X-ray spectroscopy provides information on core−valence
transitions and can probe Mn oxidation states and the ligand
environment. In previous XAS studies,15,17 conclusions were
derived only from the XANES edge region, based on
comparison with experimental results from Mn model
complexes with known structures and oxidation states. The
pre-edge region corresponds to the electric dipole forbidden 1s
→ 3d transitions, which gain intensity from metal d−p orbital
mixing when the metal ion symmetry deviates from
centrosymmetry, making this region sensitive to the ligand
field. Correlation of the pre-edge transition energies and
intensities with the structure of investigated compounds can be
performed using quantum chemistry (QM) calculations.43
However, the resolution of partial fluorescence XAS is limited
by the 1s core-hole (created upon Mn 1s excitation in XANES)
lifetime broadening, and higher resolution data are required for
rigorous and quantitative comparison with calculated spectra.
Herein, Mn Kα high-energy resolution fluorescence detected
(HERFD) spectra of the S0−S3 states are presented. The data
provide for the first time a highly resolved picture of the pre-
edge region on all S-states. Importantly, contrary to previous
studies,15,17,44 we use glycerol-free samples to avoid the
glycerol-induced nonphysiological nonphysiological heteroge-
neity in the S3 state.21,45,46 By employing time-dependent
density functional theory (TDDFT) to compute XAS spectra
of a wide variety of OEC models, we achieve a detailed
mapping between structural and spectroscopic features.
Among other important conclusions regarding the chemical
properties of distinct intermediates, correlation of the Mn XAS
pre-edge region between experiment and theory shows that the
HERFD spectra are consistent only with all-Mn(IV) oxo-
hydroxo models of the S3 state, consistent with magnetic
resonance studies. This result rules out ligand-based oxidation
during the S2 → S3 transition and locates the onset of the
formation of the O−O bond past the final light-driven
oxidation in the catalytic cycle.
■ METHODOLOGY
Photosystem II Purification and Preparation of EPR-
XAS Samples. Photosystem II was purified from the
thermophilic cyanobacterium Thermosynechococcus vestitus
(previously named Thermosynechococcus elongatus) according
to Kuhl et al.47 Based on the final purification step, highly
active PSII dimers from several preparations were pooled to
reach a homogeneous sample base. The biochemical quality of
the PSII dimers is routinely verified by Native PAGE
electrophoresis, and oxygen evolution measurements are
taken. The typical oxygen evolution activity is 5000 (±500)
μmol O2/h · mg Chl at +30 °C. Further verification of active
PSII is obtained via S-state signal progression in the EPR data.
For the present batch of samples, the standard purification
Figure 1. (a) Crystallographic structure of the OEC in the S1 state.35 (b) Catalytic cycle of water oxidation by the OEC with commonly accepted
oxidation states of the Mn ions in the S0−S2 states and the currently discussed alternatives for the S3 state.
Figure 2. Different models discussed for the S3 state, showing the
range of possible electronic structures.
Journal of the American Chemical Society
pubs.acs.org/JACS
Article
https://doi.org/10.1021/jacs.3c06046
J. Am. Chem. Soc. 2023, 145, 25579−25594
25580

\newpage

% Page 3

protocols were followed and verified by Native PAGE and
EPR, but explicit O2 activity assays were not made for the
present preparation. Given the identical EPR, however, which
evidences S state advancement, a similar activity can be
assumed. The final buffer contains 500 mM mannitol, 40 mM
MES (pH = 6.5), 10 mM CaCl2, 10 mM MgCl2, and 0.03\% v/
v n-dodecyl β-D-maltoside. Phenyl-para-benzoquinone
(PPBQ) dissolved in dimethyl sulfoxide (DMSO) was added
as electron acceptor at a final concentration of 0.5 mM.
Photosystem II (1.5 mM Mn, 15 μL) was loaded in XAS cells.
The cells were designed to fit in the EPR resonator and the
cryostat utilized for HERFD XAS measurements. Both surfaces
of the cell were open and sealed with Mylar tape windows. The
sample dimensions in the cell were 2.5 mm width and 8 mm
length and 0.8 mm depth. Samples were synchronized in the S1
state by a “pre-flash” and dark incubation at room temperature
for 1 h. All samples were illuminated by a Nd:YAG laser
(wavelength 532 nm, 10 ns pulse of 300 mJ) with 0, 1, 2, or 3
flashes and frozen in liquid N2. The laser beam was split in two
and the sample was illuminated from both sides.
EPR Quantification of S-States. The S-state advancement
is imperfect.48−56 By using short intense laser pulses, we can
avoid “double hits”, but it is impossible to avoid “misses”. The
miss factor was calculated using the intensity of the multiline
EPR signal of S2 state in all samples (Figure S1), as previously
reported by Messinger et al.15 Measurements were carried out
by using a Bruker E500 spectrometer equipped with a Bruker
ER 4116DM resonator and an Oxford Instruments ESR 935
cryostat. From the peak to peak intensity of the multiline (lines
used marked in Figure S1a), the population of each S-state
after 0, 1, 2, and 3 flashes was calculated (Tables S1, S2) and
used in order to deconvolute the pure S1, S2, S3, and S0 XAS
spectra. We note that after the fourth flash the intensity of the
S2 multiline signal is 80\% compared to the first cycle (Figure
S1b). The very good conversion factor confirms the high
activity of the samples. Given that inactive OEC centers would
have released Mn(II) in solution, all samples were checked for
the presence of Mn(II) EPR signals and batches with high
Mn(II) content were discarded. Crucially, the homogeneity of
our S3 state samples has been confirmed by W-band EPR
spectroscopy (Figure S2).45
HERFD Measurements. Mn Kα HERFD measurements
were performed at beamline 6−2 of the Stanford Synchrotron
Radiation Lightsource (SSRL).57 The SPEAR storage ring was
operating at 3 GeV with a current of 500 mA. BL6−2 utilizes a
56-pole 0.9 T wiggler insertion device. Two Rh-coated Si
crystals (one flat for vertical collimating and one cylindrical for
focusing, upstream and downstream of the monochromator)
were utilized to focus the beam. The incident beam was
monochromatized using a pair of cryogenically cooled Si(311)
crystals, giving an energy resolution of ∼0.2 eV. The flux was
estimated to be ∼2.9 × 1011 photons/s in a 420 × 200 μm
beam spot. The energy of the incident beam was calibrated by
setting the maximum of the pre-edge peak of KMnO4 to
6543.21 eV. Kα-HERFD and Kα XES data were collected by
utilizing a Johann-type XES spectrometer equipped with five
Ge(111) crystals and an energy-dispersive silicon drift
detector. The spectrometer resolution at the elastic peak was
0.8 eV. The sample temperature during measurements was
poised at 10 K by using liquid He-cooled flow cryostat. The
beam shape was Gaussian with vertical full width at half height
of 200 μm and horizontal of 420 μm. By using a vertical step of
300 μm and horizontal step of 630 μm, it was possible to
measure ∼50 separate spots per sample. The Kα HERFD scans
were collected in an energy range of 6530−6570 eV, while
longer 6530−6800 eV scans were measured to facilitate proper
normalization. The energy step was varied across the
measurement and was 0.15 eV for the pre-edge region,
which is the primary focus of the present study, and 0.2 eV
across the edge. The Kα XES spectra were collected to assess
the Kα maxima of each S state, and hence, short scans were
used in the energy range of 5896−5904 eV.
Damage Studies. Detailed damage studies were per-
formed before the actual measurement in all S states (Figure
S3). They involve measurement of XAS scans in the range
6535−6570 eV in order to establish the maximum dwell time
per sample spot. Due to the low sample concentrations, we
exposed to radiation a number of spots and collected
consecutive scans in each spot. We compared the average of
all the first scans, to the average of all the second scans and so
on, to establish the time point at which a change in the pre-
edge and edge region is observed. Attenuation of the incident
beam by using foils was also required. A range of attenuation
factors was tested, and we concluded that 6\% of the total flux
was not damaging the sample during a scan time of ∼60s, while
reasonable signal/noise was maintained. Additionally, time
scans were performed at a constant incident energy of 6551 eV,
the energy point that was considered as the most sensitive
probe of damage during the initial damage scans. The dwell
time for S1 and S2 states was determined as ∼100 s/spot while
for S3 and S0 was ∼75 s/spot and was used during the actual
measurement. Higher oxidation states of the Mn4Ca cluster are
expected to damage faster. This fact was verified by the damage
studies, and a shorter scan (75 s vs 100 s/spot) was used for
the S3-state. We note that while a shorter dwell time for the
formally more reduced S0 may seem counterintuitive, this is
due to the fact that the S0 samples as prepared contain a
significant amount of S3 as quantified by EPR. A dose of 1.7 ×
107 photons/μm2 was used for the S1 state, 1.4 × 107 photons/
μm2 for the S2 state and 1.2 × 107 photons/μm2 for the S3 and
S0 states.
HERFD Data Processing. Every scan was normalized to
the incident flux I0 and subsequently scans of the same flash
number were averaged in PyMCA in order to create the
average of 0, 1, 2, and 3 flash spectra (ca. 150 scans were
averaged for the S1 and S2, while ca. 220−230 scans for the S3
and S0 states). All other processing was performed in Matlab.
The short scans (6530−6570 eV) were normalized to the edge
jump by using the long scans (6530−6800 eV). Spectra of pure
Si states were acquired by subtraction of the appropriate
percentage of pure Si−1 (and Si−2 if needed) as calculated by
EPR quantification of S-states after each flash (Tables S1 and
S2). The rising edge of the pure spectra was fitted by
interpolating a spline polynomial in the region 6530−6550 eV.
By subtraction of the XAS spectrum minus the fitted line, the
isolated pre-edge spectra of the S-states were acquired. Fitting
of the pre-edge peaks was performed using Voigt curves
(Lorentzian has only a small contribution) and least-squares
fitting. The fact that the peaks represent Gaussian shape is
expected since the 1s core hole lifetime broadening
(Lorentzian) is effectively eliminated in a HERFD experiment,
and the energy resolution is thus dominated by the
instrumental broadening (Gaussian).
Evaluation of the Error. The error of the 0, 1, 2, and 3
flash spectra was estimated by the standard error of the mean
of averaged spectra (Figure S4). This is a measure of how far
Journal of the American Chemical Society
pubs.acs.org/JACS
Article
https://doi.org/10.1021/jacs.3c06046
J. Am. Chem. Soc. 2023, 145, 25579−25594
25581

\newpage

% Page 4

our sample mean, i.e., scans of our experiment, is likely to be
from the true population mean that would be infinite scans
that would give the real spectrum. The standard error is
calculated as
x =
n where
=
= x
x
n
(
)
1
j
n
j
1
2
is the standard
deviation, x̅ is the mean of scans x1−xn, and n is the total
number of scans. After subtraction of the appropriate
percentage of the Si−1 state determined by EPR from the raw
spectra in order to acquire the pure Si state spectra and
renormalization,
the
error
was
recalculated
as
=
+
X
X
X
X
X
S
S
S
S
1
1
2
2
2
2
Si
Si
Si
i
i
i
i
1
1
1
1 , where
X
X
Si
Si
1 is
the standard error of the pure spectrum of the Si state, XS
di is the
fraction of the Si state in the sample, while XS
di−1 is the fraction
of the “missed” centers in the Si−1 state that was subtracted, σS
di
is the standard error of the spectrum of Si, and σS
di−1 is the
standard error of the Si−1 spectrum.
Computational Details. Cluster models of the OEC were
constructed from the X-ray crystallographic coordinates of the
5B66 structure monomer B.35 The models include the
inorganic core Mn4CaO5, terminal water molecules W1−W4,
and first coordination sphere amino acids His332, Glu189,
Asp342, Ala344, CP43-Glu354, Asp170, and Glu333. Second
coordination sphere residues include His337, CP43-Arg357,
Asp61, Ser169, Leu343, Tyr161, His190, Asn298, Gln165,
Val185, and Phe186, and 13 crystallographic water molecules:
6 from the O1 water channel, 3 hydrogen bonding to W3 of
the Ca2+ ion, 2 from the Cl− water channel, and 2 from the O4
water channel. Overall, the S1 state model (with W2 = H2O)
includes 329 atoms and is shown in Figure S5.
All calculations were performed with Orca 5.58 Geometry
optimizations were carried out using the B3LYP59,60 functional
with the resolution of the identity (RI) approximation
including D4 dispersion corrections.61 Optimizations were
performed in the respective broken symmetry states for each
model, i.e. α−β−α−β for S1A, S1HA, and S0O5, β−α-α−α for
S1B, S1HB, S3P, and S0HO4, α−β−β−α for S1O5H, S2A,
S2HA, and S0O4, α−α−α−β for S2B, S2HB, S2O4H, S3AW,
and S3BW, α−β−α−α for S0HO4, and α−α−α-α−β (O5−O6
radical antiferromagnetically coupled to Mn ions) for S3O.
Relativistic effects were considered throughout using the
zeroth-order regular approximation (ZORA).62−64 The
scalar-relativistically recontracted65 ZORA-def2-TZVP(-f)
basis sets66 were used for all atoms except C and H for
which the ZORA-def2-SVP basis sets were used. In all
calculations, the CPCM solvation67 with dielectric constant ε
= 6.0 was used to simulate the effect of the protein
surrounding.
XAS spectra were calculated with the TD-DFT method
employing the Tamm−Dancoff approximation.68,69 The
TPSSh functional70 was used and the computed excitation
energies were shifted to higher energies by 36.3 eV, to account
for systematic methodological deviations, based on previous
benchmarking studies on Mn monomers and dimers.71,72 The
accuracy of the method has been quantified in previous
benchmark studies on mononuclear Mn complexes using the
computed R2 of the linear relationship between the calculated
and experimentally determined transition energies (R2 = 0.94)
and intensities (R2 = 0.88).72 The RI and chain of spheres
(RIJCOSX) approximations73 were employed to speed up the
calculations. The XAS spectra of each Mn ion were calculated
with 150 roots, and the spectra of each model were plotted by
adding the spectra of the four Mn ions and using a peak
Gaussian broadening of 1.1 eV. In order to compare with the
normalized intensities of the experimental spectra, all
computed intensities were multiplied by 0.02, based on
agreement between the maximum intensity of the 0F state
with the computed intensities of the S1 states.
■ RESULTS AND DISCUSSION
HERFD Experiment. The Mn Kα HERFD X-ray
absorption spectra of the dark-adapted S1 state, the one-flash
(S2), two-flash (S3), and three-flash (S0) illuminated states of
PSII samples from T. vestitus were measured at beamline 6−2
of the Stanford Synchrotron Radiation Lightsource (SSRL)
(Methods section, SI) and are shown in Figure 3a. All spectra
are normalized to the edge jump, as described in the Methods.
Functionality of the samples is indicated by the period four
oscillation in the S2 state multiline g ≈ 2 EPR signal (Figure
S1).
Progression of PSII samples to the next S-state upon
illumination is not perfect; thus, quantification of the
population of PS II centers of the sample in each S-state is
essential. EPR spectroscopy was employed following a
previously suggested protocol.15 Quantification of the flash-
induced S-state progression efficiency was performed based on
the intensity of the multiline g ∼ 2 EPR signal of the S2 state
after each flash.15 The S2 population after each visible light
flash was quantified, and a model considering misses for every
Figure 3. (a) HERFD spectra of the dark-adapted samples (0F) and after 1−3 flashes (1−3F). Spectra are normalized to the edge jump. 3-point
adjacent averaging was used. (b) Pure spectra of all S-states after the appropriate subtractions based on EPR quantification.
Journal of the American Chemical Society
pubs.acs.org/JACS
Article
https://doi.org/10.1021/jacs.3c06046
J. Am. Chem. Soc. 2023, 145, 25579−25594
25582

\newpage

% Page 5

flash illumination was fitted to the EPR data. The miss factor
was calculated to be 6\%. In other studies, miss factors of 6−
18\% have been reported.10,15,17,19,44 In a separate experiment, a
singly flashed sample was illuminated continuously at 200 K in
order to obtain the maximum of S2 and no increase in the
multiline signal intensity was observed. Double hits did not
improve the fitting. A considerable population of PS II centers
(15\%) remains in the S2 state and does not progress to S3 and
S0 upon illumination. This is caused by acceptor side
deficiency: the acceptor side in some PSII centers cannot
accept more electrons in order for the OEC to proceed further
than S2. Single turnover values between ∼5 and 10\% of PSII
centers have been reported elsewhere.10,15 A percentage
(∼5\%) of centers in S2 state in the dark adapted sample has
also been considered, but the inclusion of such population did
not improve the fitting (see also SI, section 1 and Table S1).
The population of each S-state after each laser pulse is
presented in Table 1. After “0, 1, 2, 3 flashes”, there is 100\% S1,
94\% S2, 74\% S3, and 70\% S0, respectively.
The optimal fitting with the experimental data was achieved
by a miss factor of 6\% and assuming that a 15\% population of
the OEC centers remains blocked in the S2 state. The pure S-
state spectra derived after deconvolution according to the S-
state quantification are shown in Figure 3b. The spectra consist
of the rising edge at 6545−6560 eV, that represent 1s → 4p
transitions, and the pre-edge at 6538−6545 eV, attributed to
dipole forbidden 1s → 3d transitions. Each of the regions is
examined separately in the next sections.
Analysis of the Edge Region. For systems with similar
coordination environment, the energy of the edge is dependent
on the oxidation state of the metal ion: as the oxidation state
increases the edge shifts toward higher energy. Visual
inspection of the edge regions in Figure 3b shows that the
edge shifts to higher energies upon the S0 → S1 and S1 → S2
transitions. For the S1 → S2 transition, the edge shift can be
attributed only to the more positive charge of the cluster since
the structures of S1 and S2 are similar, as proposed by
EXAFS17,44 and most recently by the XFEL crystal
structures.29 The S0 → S1 transition involves deprotonation
in addition to oxidation; therefore, the observed edge shift is
attributed to Mn-centered oxidation as well as small structural
rearrangements. Thus, in agreement with previous
works,15,17,44 positive edge energy shifts are observed during
the S0 → S1 and S1 → S2 transitions, which reflect Mn-based
oxidations.
The edge shift during the S2 → S3 transition is not equally
obvious from visual inspection because the shapes of the edge
of S2 and S3 states are significantly different, as shown also in
the first derivatives of the spectra (Figure S6). Those
differences possibly stem from significant structural changes
involved in the S2 → S3 transition, which include
deprotonation and possibly the insertion of a water molecule
in the cluster. Thus, in S2 → S3, the edge shift is dependent on
the oxidation state changes as well as the ligand field change;
since the magnitudes of these effects are comparable, this
complicates the interpretation of the edge as it has been shown
in Mn model systems.74 A change in the coordination
environment of a Mn ion from 5-coordinate Mn(III) to 6-
coordinate Mn(IV) was proposed based on comparison of Kβ
emission difference spectra of the S3 minus the S2 state with
the difference spectra of synthetic models with Mn(III)L5,
Mn(IV)L5, and Mn(IV)L6.9
Three different methods were used for assessing the energy
of the edge in XAS experiments: (a) the position at the half
intensity of the normalized spectra,75 (b) the integral
method,76 and (c) the energy at the inflection point.15 In
Table 2, the energies of the edge and the edge shifts
determined using all of the above methods are given. The
half intensity energy point, after proper normalization using the
long scans up to 6800 eV, shifts ∼0.7 eV with each S-state
transition from S0 to S3. Similar results are obtained with the
integral method, which gives a shift of ∼0.6 eV. By contrast,
the inflection point of the edge, determined by the second
derivative (Figure S7), is almost the same for the S2 and S3
states. Notably, the present results for the inflection point
values are similar to those reported by Messinger et al.,15 while
those reported with the integral method are consistent with
Haumann et al.17 As shown in Figures S8 and S9, the present
XANES spectra are very similar in the rising edge region to the
previously reported XANES measured on spinach15,17 and
cyanobacteria.44 The inflection point method is more sensitive
to the shape of the edge since it is defined by a single point.
Thus, the shift of the inflection point energy is not consistent
with the other two methods and the different methods used by
the different groups for assessing the edge shift may cause the
discrepancy in the literature regarding the presence17,75,77 or
near absence15,78 of edge shift during S2 → S3 and thus if Mn-
centered or ligand-based oxidation occurs during this
transition.12,22,32,38,41,79−82
Overall, all methods of determining the edge energy confirm
Mn-based oxidation in the S0 → S1 and S1 → S2 transitions. On
the other hand, diverging results are obtained in the case of the
S2 → S3 transition by different methods, which has given rise to
seemingly contradictory interpretations among different
groups, despite their data being almost identical. Given the
additional complication of structural differences between the S2
and S3 states, the locus of oxidation in this transition cannot be
conclusively determined from the analysis of the edge region.
Therefore, in the next sections we focus on the analysis of the
pre-edge region, which more directly correlates with the
electronic structure, and where the high resolution of the
HERFD spectra combined with QM calculations provide an
unambiguous interpretation.
Analysis of the Pre-edge Region. Figure 4 shows the
expanded pre-edge region for all S-states. Before analyzing the
Table 1. S-State Content (\%) after a Given Number of
Flashes
flash
S1
S2
S3
S0
0
100
1
6
94
2
1
25
74
3
0
16
14
70
Table 2. Edge Energies and Edge Shifts (Si−Si‑1, in
Parentheses) Obtained with Three Different Methodsa
half of normalized Intensity
integral method
inflection point
S0
6550.8
6551.6
6550.5
S1
6551.4 (0.6)
6552.1 (0.5)
6553.1 (2.6)
S2
6552.2 (0.8)
6552.7 (0.6)
6553.7 (0.6)
S3
6552.9 (0.7)
6553.2 (0.5)
6553.9 (0.2)
aAll energies in eV.
Journal of the American Chemical Society
pubs.acs.org/JACS
Article
https://doi.org/10.1021/jacs.3c06046
J. Am. Chem. Soc. 2023, 145, 25579−25594
25583

\newpage

% Page 6

current data in more detail, it is important to compare these
spectra to previous XAS data. As discussed above, the XANES
region of the HERFD spectra are essentially identical to
previous reports.15,17,44 However, as seen in Figure S8, the
ratio of intensities in the pre-edge region is modulated relative
to previous reports. First, we note that the present data
represent the highest resolution XANES data available for all S
states. While 1s2p RIXS planes of the S-states were previously
reported by Glatzel et al.,83 these data do not include the full
XANES and thus do not allow for proper normalization.
Further, the previous data were collected with a Si(111)
monochromator relative to the Si(311) monochromator used
in the current study. Hence, we estimate the previous report
would have a resolution of ∼1.3 eV at constant emission
energy relative to the ∼0.8 eV resolution in the current study.
These differences in resolution, however, are unlikely to
account for the differences in the pre-edge ratios. More likely,
the observed differences are attributed to the fact that no
glycerol was used in the present study. As glycerol is known to
change the equilibrium of the different configurations of the S
states, we hypothesize that this may be the origin of the
modulations in the pre-edges. This is consistent with previous
studies by Yano et al.,46 which showed that near-infrared
conversion of the S2 low spin multiline signal to a high-spin
EPR signal is correlated with a decrease in the intensity of the
first pre-edge feature. These observations thus further motivate
the importance of having the present data set for all S-states in
the absence of glycerol. In order to quantitatively correlate the
present data with theory, we proceeded to a detailed fitting
analysis.
The pre-edge region of each state was fitted by Gaussian
curves, and the optimal fittings are shown in Figure 4a. The
pre-edge regions of the S0 and S1 states can be best described
by three features: a lower intensity peak at 6540 eV and two
higher intensity peaks at 6541 and 6543 eV. The S2 and S3 pre-
edge were fitted by only two peaks at 6541 and 6543 eV.
Alternative fittings are given in Figure S10 and Table S3.
Changes in the pre-edge region with S-state progression can
be quantified based on the areas of the individual peaks and of
the overall pre-edge region, given in Table 3. The intensity of
the first peak decreases from the S0 to the S1 peak, and the
peak is not detectable at all at S2 and S3. The second peak at
6541 eV becomes more intense in S1 than in S0 and has a
similar intensity in S1−S3. A subtle peak is also observed at
∼6546 eV in the S0 spectrum partially covered by the rising
edge. The third peak at 6543 eV increases from S0 to S1 and
from S2 to S3, but not in S1 to S2. In order to compare solely
the contributions to the XANES spectra due to the 1s → 3d
Figure 4. (a) Fitting of the pre-edge region with Gaussian peaks and the rising edge background. (b) Pre-edge region of each S-state after the
subtraction of the rising edge background. (c) S-state pre-edge XAS spectral differences. The dashed lines show the zero value for each difference
spectrum.
Table 3. Energy in eV and Area (in parentheses) for Each Fitted Peak in Each S-State Pre-edge Spectrum and Intensity
Weighted Average Energy (IWAE) and Area of the Total Pre-edge Region
peak 1
peak 2
peak 3
total
S0
6539.9(0.031)
6540.9(0.082)
6542.5(0.079)
6541.4(0.208)
S1
6539.7(0.008)
6540.9(0.118)
6542.7(0.115)
6541.6(0.235)
S2
6540.8(0.108)
6542.8(0.108)
6541.8(0.217)
S3
6541.0(0.113)
6542.9(0.121)
6542.0(0.238)
Journal of the American Chemical Society
pubs.acs.org/JACS
Article
https://doi.org/10.1021/jacs.3c06046
J. Am. Chem. Soc. 2023, 145, 25579−25594
25584

\newpage

% Page 7

transitions, the contribution of the rising edge to the pre-edge
was subtracted. The spectra of S0−S3 with the subtracted
baseline are shown in Figure 4b. The pre-edge region changes
with S-state progression are shown in the difference spectra in
Figure 4c. Figure S11 indicates the standard errors of the
difference spectra.
The changes in the overall shape of the pre-edge region
during the S0−S3 state transitions can be expressed through the
intensity weighted average energies (IWAEs). The IWEAs shift
0.2 eV higher after each transition, which imply that higher
energy transitions gain intensity with each Mn oxidation.
Notably, those changes cannot be explained solely in terms of
Mn oxidation during each transition, as in the case of the edge
region. Apart from indirect correlations between pre-edge
intensities and the presence of specific groups,84 a detailed
correlation between structure and spectroscopy requires QM
calculations.71,72,84−87 Therefore, in the next sections, we
attempt to correlate the pre-edge features with structural
features using QM-derived models of all states.
Correlation with Structural Models. To correlate the
XAS spectra with the structure of the OEC in each state, we
optimized models of several variants of the S0−S3 states, and
we calculated their XAS spectra. The core parts of all
optimized structures are shown in Figure 5, along with
explanations of the labeling used in this work. Additionally, the
inorganic cores of all models are depicted schematically in
Figure S12 where the Jahn−Teller axis orientations for all
Mn(III) ions are depicted and important distances are
indicated. For the S1 state, we considered five models with
variations in the protonation states of terminal water-derived
ligands on Mn4 and the O5 bridge, as well as in the orientation
of pseudo-Jahn−Teller elongation axes of the Mn(III) ions.88
In all optimized S1 state models, the Mn oxidation states are
III−IV−IV-III. For the S2 state, we considered five models that
take into account all different ideas discussed in the literature
regarding protonation states of terminal water-derived ligands
and the O4 bridge,89 as well as valence isomerism (i.e., Mn
oxidation states III−IV−IV−IV and IV−IV−IV−III).90 For
the S3 state, we considered all possibilities discussed in the
introduction regarding Mn and ligand oxidation states, i.e., Mn
IV−IV−IV−IV with oxo and hydroxyl O5 and O6 ligands,
III−IV−IV−IV with oxo-oxyl radical, or III−IV−IV−III with
peroxo O−O bond formation, in total 5 models. Finally, for
the S0 state we considered 4 models differing in the
protonation states of W2 and of the O4 and O5 bridges.91,92
The TD-DFT calculated XAS spectra of all 19 optimized
structures are shown in Figure S13, compared to the
experimental spectra for each S-state.
Analysis of Difference Spectra. The difference spectra
between successive S-states are substantially more informative
compared to the individual spectra, as it removes systematic
errors in the intensity evaluation and explicitly isolates the
features that change from one state to the next. The results
show that in some cases where the similarities in computed
spectra are significant, the difference spectra enable the
confident selection of the best fitting structural models.
Figure 6 compares the S2−S1 difference spectra between
experiment and theory for the selected computational models.
Importantly, the comparison of difference spectra using
distinct computational models is so sensitive that it is possible
to distinguish between S1 state isomers in which the
orientation of the Jahn−Teller axis of the Mn4(III) ion is
either perpendicular or collinear to the Mn1(III) Jahn−Teller
axis (S1
A,H and S1
B,H, respectively),88 but also between different
protonation states of the terminal W2 ligand in both the S1 and
S2-state models. For example, the S2−S1 difference spectra for
Figure 5. Optimized core structures of the computational models of S0−S3 states considered in this work. For clarity, only the inorganic core and a
few selected coordinating ligands (of the actual ∼330 atom models) are depicted in this figure. Mn(III) ions are shown in light pink, and Mn(IV)
ions are in dark purple. Label superscripts A and B are for “open” and “closed” cubane geometries, respectively; H superscript denotes that W2 is in
the aquo form (H2O), otherwise, it is in the hydroxo form, and O4 and O5 denote that bridging oxo ligands O4 and O5, respectively, are
protonated in the specific models. The models whose labels are in colored boxes turned out to be in best agreement with experiment, as explained
in the following sections.
Journal of the American Chemical Society
pubs.acs.org/JACS
Article
https://doi.org/10.1021/jacs.3c06046
J. Am. Chem. Soc. 2023, 145, 25579−25594
25585

\newpage

% Page 8

the S1 models with perpendicular Jahn−Teller axes (models
S1
A and S1
A,H) show a large negative intensity difference at
6542−6543 eV, inconsistent with experiments, whereas this is
well reproduced with model S1
B,H. Interestingly, the observed
negative intensity difference in the region 6539−6540 eV is
only reproduced by S1 models that have W2 in the aquo form
(in agreement with FTIR studies93), but not by those with W2
in the hydroxo form. This is associated with the shoulder at
6539.5 eV that is observed for the S1 models where W2 is H2O
(Figure S13a).
The models that best reproduce the experimental S2−S1
difference spectra are S1
B,H and S2
A,H (Figure 6a, blue line).
The comparison shows the excellent agreement between the
experiment and calculations. It is worth pointing out that the
calculated XAS spectrum of S2
B,H is slightly lower in intensity
than S2
A,H (Figure S13b), similar to the S2 high-spin and low-
spin forms reported by Chatterjee et al.46 By contrast, the
computed spectrum of the S2 state model with protonated O4,
(S2
O4H), which has been suggested as a possible interpretation
of the high-spin S2 form,89,94 is more intense than S2
A,H (Figure
S13b). Therefore, our experimental pre-edge results combined
with TD-DFT calculations are most consistent with valence
isomerism90 giving rise to the low and high spin S2 state
signals, respectively. As the samples were not treated with
glycerol,95 about 20\% of the centers are in the high-spin S2
state (quantification based on EPR, see SI). Since the fraction
of high-spin S2 is small, the effect on the S2−S1 difference
spectra in the pre-edge region is negligible (Figure S14) and
thus we use only the dominant S2
A,H spectrum for the
difference spectra plots.
Comparison of the calculated S1−S0 difference spectra using
different computational models, in Figure S15, shows that
model S0
O4 exhibits the best fitting with experiment among the
different S0 variants. Notably, light-induced Fourier transform
infrared (FTIR) S1−S0 difference spectra reported by
Yamamoto et al.96 also indicate O4 protonation in the S0
state; however, that study also supports a fully protonated W2,
which would be represented by our S0
O4,H model. This would
indeed be more consistent with one electron oxidation and one
deprotonation (of O4) during the S0 → S1 transition, and it
needs to be further investigated. Previous benchmarking
studies on Mn oxo-bridged compounds have also demon-
strated the sensitivity of Mn XAS to oxo-ligand protonation
states.71 Therefore, analysis of the S0−S2 state transitions
shows that the pre-edge region combined with quantum
chemistry simulations not simply confirms Mn-based oxidation
but directly constrains the possible structural features of each
state.
Even more pronounced than the S1 and S2 states are the
differences observed among computational models of the S3
state in the case of S3−S2 difference spectra shown in Figure 7.
The comparison reveals that the Mn(IV)4 oxo-hydroxo models
(S3
A,W and S3
B,W) are clearly in better agreement with
experiment than the Mn(III)Mn(IV)3 oxyl-oxo (S3
oxyl−oxo)
and Mn(III)2Mn(IV)2 peroxo (S3
peroxo) models. The S3
oxyl−oxo
and S3
peroxo are completely inconsistent because S3
oxyl−oxo
shows much higher intensity than S2 between 6541 and
6544 eV, whereas S3
peroxo shows much lower intensity. Among
the Mn(IV)4 oxo-hydroxo structural isomers, S3
A,W and S3
B,W,
the open cubane S3
A,W shows the closest agreement with
experiment. The above results strongly disfavor ligand-based
oxidation during the S2 to S3 transition.
Analysis of S-State Progression. In Figure 8, we show
the calculated spectra of the most experimentally consistent
models of each state, as identified from the analysis of the
difference spectra (structures S1
B,H, S2
A,H, S3
A,W, and S0
O4, as
labeled in Figure 5). The most experimentally consistent
Figure 6. Experimental and calculated S2−S1 spectral differences for
different S1 and S2 isomers.
Figure 7. Experimental S3−S2 data and calculated S3−S2 difference
spectra for models S3
A, W, S3
B, WS3
oxyl−oxo, S3
peroxo, and S2
A, H.
Figure 8. Calculated Mn K-edge XAS spectra for the most
experimentally consistent models of states S0 to S3.
Journal of the American Chemical Society
pubs.acs.org/JACS
Article
https://doi.org/10.1021/jacs.3c06046
J. Am. Chem. Soc. 2023, 145, 25579−25594
25586

\newpage

% Page 9

model of the S3 state (S3
A,W) is an all-Mn(IV) model, and
hence, the XAS pre-edge region is consistent with Mn-based
oxidation in each S-state transition. Importantly, the same
model is the one that fully explains observations by electron−
electron double resonance (ELDOR) detected nuclear
magnetic resonance experiments (EDNMR).20,21
The calculated spectra in Figure 5 reproduce all of the major
features of the experimental spectra. Specifically, (a) the
shoulder at 6539.5 eV is smaller in the S1 state than in the S0
state and is absent in the S2 and S3 states, as in experiment; (b)
the peak above 6541 eV has the smallest intensity in the S0
state and the largest in the S3 state; (c) the intensity in the
region between 6542 and 6544 eV increases with each S-state
transition; and (d) the intersection points of all spectra at
6540.5 and 6544 eV, and of the S1 and S2 states spectra at
6542.7 eV are also reproduced. Notably, small deviations from
experiments in the relative intensities among the S-states are
within the accuracy limits of our methodology and can be
more easily distinguished in the corresponding difference
spectra in Figures 6, 7 and Figure S15.
It is important to note that direct visual comparison of the
calculated spectra of a specific state with experimental spectra
is not expected to lead to useful conclusions; instead, it is
meaningful to compare the total pre-edge area as well as the
IWAEs to experiment, as stressed in previous studies.71,72,97
Each spectrum results from the distribution of energies and
intensities of a vast number of transitions from different Mn
ions, and the height and width of each apparent peak of the
spectrum are very sensitive to small variations in the relative
energies of the transitions it contains. Moreover, it is amply
established that errors in the (relative) intensities of the
transitions are larger than errors in the energies.72 It is
precisely for these reasons that the pre-edge area is considered
a far better parameter than individual peak heights to assess
agreement with experiment.97 Indeed, an excellent correlation
is observed between summed calculated intensities and
experimental pre-edge areas, as presented in Figure 9a. The
linear relationship (R2 = 0.98) shows a strong correspondence
between experiment and theory. It also implies that intensity
errors are similar among different S-states, which further
justifies the use of difference spectra to probe S-state
progression. By contrast, the transition intensities of the
ligand-based oxidized alternatives for the S3 state, S3
oxyl−oxo and
S3
peroxo (Figure 2), deviate strongly. The computed spectra also
reproduce the shift of the IWAEs to higher energy with the S-
state progression (Figure 9b).
Electronic Structure Origin of the Observed Spectral
Features. The contributions of each Mn ion to the calculated
XAS spectra of the most experimentally consistent models of
each S-state are shown in Figure 10, where the stick spectra of
each Mn ion are plotted in different colors for each structure.
Analysis of the TD-DFT natural transition orbitals (NTOs)
reveals the nature of the underlying transitions that form the
pre-edge region. The NTOs of S1
B,H, S2
A,H, S3
A,W, and S0
O4 as
well as of the redox isomers of the S3 state, S3
A,W, S3
oxyl−oxo, and
S3
peroxo, are shown in Figures S16−S21. We observe that the
nature of the transitions in each energy region remains
essentially the same with S-state progression; thus, in Figure 10
we focus on the origins of the differences between successive S-
states, and we show the NTOs that contribute the most to
those differences.
Specifically, transitions with energy lower than ∼6542 eV are
mostly attributed to excitations from the 1s core orbitals to 3d
with dominant local Mn character, while transitions with
energy higher than 6542 eV also have metal-to-metal charge
transfer character. Transitions from the 1s to nonbonding t2g
orbitals have very low intensity due to negligible metal 4p
mixing, while transitions to the antibonding dz
2 and dx
2−y
2
orbitals contribute significantly to the intensity due to
increased metal 4p mixing, which is covalently mediated by
σ interacting oxygen 2p orbitals. The most intense spectral
differences are observed in regions A−D, denoted with pink,
purple, and blue hues in the stick spectra plots in Figure 10.
In S0
O4, the peak around 6540 eV (region A) is mostly
attributed to local Mn4(III) and Mn3(III) 1s → 3d transitions,
whereas in S1
B,H the latter is blue-shifted due to Mn3(III)
oxidation to Mn3(IV), and it appears in region B. In addition,
Mn3(IV) transitions in regions B, C, and D gain intensity in
S1
B,H, which explains the positive S1−S0 difference spectra in
those regions (Figure S15).
As in the S0 → S1 transition, Mn4(III) oxidation to
Mn4(IV) during the S1 → S2 transition blue-shifts the local
Mn4 1s → 3d transition from region A in S1
B,H to region B in
S2
A,H, and local Mn4(IV) transitions gain intensity in regions B
and C. However, the local Mn3(IV) transition at 6541.5 eV
has decreased intensity in S2
A,H relative to S1
B,H, probably due
to decreased Mn 3d−4p orbital mixing as a result of
coordination geometry changes, leading to negative S2−S1
difference spectra in region B. Furthermore, a metal-to-metal
charge transfer Mn3(IV) 1s → Mn4(IV) 3d transition arises in
region D. This qualitative analysis reveals that the effect of
metal-based oxidation on the XAS pre-edge region of the OEC
Figure 9. (a) Correlation of summed calculated intensities and
experimental pre-edge areas and (b) correlation of calculated and
experimental intensity-weighted average energies for different OEC
models in different S-states. The correlations show that only Mn-
centered oxidation in the S2 → S3 transition is consistent with the
HERFD data.
Journal of the American Chemical Society
pubs.acs.org/JACS
Article
https://doi.org/10.1021/jacs.3c06046
J. Am. Chem. Soc. 2023, 145, 25579−25594
25587

\newpage

% Page 10

is essentially due to increased intensity of transitions that
involve the oxidized Mn(IV) ion, consistently with previous
reports,97 as well as due to coordination geometry changes on
the rest of the Mn ions of the cluster.
During the S2 → S3 transition, Mn1(III) oxidation to
Mn1(IV) blue shifts local Mn1 transitions, leading to negative
S3−S2 difference spectra at 6540−6541 eV. Mn1(IV)
transitions are more intense in S3
A,W than S2
A,H in regions B,
C and D, whereas Mn4(IV) transitions are less intense. The
effect of those changes is reflected in the calculated S3−S2
difference spectra, which is in excellent agreement with
experiment (Figure 7). Interestingly, the predicted pre-edge
regions of the alternative S3-state models S3
oxyl−oxo and S3
peroxo
are strikingly different.
The stick spectra of S3 state isomers S3
A,W, S3
oxyl−oxo, and
S3
peroxo are compared in Figure 11. Comparison of the stick
spectra reveals that the large intensity difference between
S3
oxyl−oxo and S3
A,W peaks is attributed to regions B, C, and D.
Notably, region A has a local Mn4(III) character. The NTOs
that correspond to peaks B−D are shown next to each plot in
Figure 11. NTOs that correspond to excitations from Mn1 1s
core orbitals are shown in green boxes and from Mn4 in the
orange box. In addition, the calculated spectra of individual Mn
ions are compared in Figure S22. The higher intensity of the
transitions of the S3
oxyl−oxo model compared to the S3
A,W in the
6541−6544 eV region is attributed mostly to excitations from
Mn1 and Mn4 1s core orbitals (Figure S22). Comparison of
the NTOs of the B, C and D peaks of S3
oxyl−oxo and S3
A,W
shows that (a) the local 1s to 3d transitions of Mn1 (first row
Figure 10. Assignment of the calculated pre-edge XAS spectrum based on the NTOs associated with the transitions for models S0
O4, S1
B,H, S2
A,H,
and S3
A,W. In the stick spectra, transitions from 1s core orbitals of Mn1 are shown in green sticks, of Mn2 in red, of Mn3 in blue, and of Mn4 in
orange. NTOs for transitions from 1s core orbitals of Mn1 are shown in green boxes, of Mn2 in red, of Mn3 in blue, and of Mn4 in the orange box.
Journal of the American Chemical Society
pubs.acs.org/JACS
Article
https://doi.org/10.1021/jacs.3c06046
J. Am. Chem. Soc. 2023, 145, 25579−25594
25588

\newpage

% Page 11

of S3
oxyl−oxo NTOs in Figure 11) are more dipole-allowed due
to the increased p character of the acceptor orbitals and are
thus more intense than the corresponding S3
A,W peaks, and (b)
S3
oxyl−oxo additionally has three intense charge transfer
excitations from Mn1 and Mn4 to O5 and O6 p orbitals
(second row of S3
oxyl−oxo NTOs). Thus, examination of the
NTOs of S3
A,W and S3
oxyl−oxo shows that the higher calculated
intensity for the S3
oxyl−oxo model is attributed to charge transfer
channels enabled by the oxyl-oxo radical.
The lower intensity of the transitions in S3
peroxo compared to
S3
A,W in the 6541−6544 eV region is traced to the same B, C,
and D peaks of S3
A,W as in the comparison between S3
A,W and
S3
oxyl−oxo. The corresponding peaks in S3
peroxo have lower
intensity, in the case of B and C, or do not exist at all, in the
case of D. In the S3
peroxo acceptor, orbitals of local excitations
of Mn1(III) and Mn3(IV) have less p character than the
corresponding excitation in S3
A,W, indicating that the peroxo
unit hinders charge transfer channels, contrary to the oxyl-oxo
and oxo-hydroxo groups.
In summary, a comparison of the experimental and
calculated pre-edge XAS spectra supports Mn(III) to Mn(IV)
oxidation in each S-state transition. Ligand-based oxidation in
the S2 → S3 transition would result in either significantly more
intense (oxyl formation) or significantly less intense (peroxo
formation) transitions, therefore this possibility can be safely
excluded based on the present HERFD data.
■ CONCLUSIONS
We have presented HERFD spectra of all S-states of PSII from
T. vestitus. Spectra were collected as “0, 1, 2, and 3 flashed”
samples and deconvoluted to pure S1, S2, S3 and S0,
respectively, with the aid of EPR spectroscopy. The energy
of the edge shifts to a higher energy during the S0 to S1 and S1
to S2 transitions, which clearly reflects Mn-based oxidation.
However, the edge shift cannot be used reliably to determine
the locus of oxidation during the S2 to S3 transition, due to
inherent complexity of this oxidation reaction as well as due to
quantification method limitations. Therefore, we focus on the
pre-edge region, leveraging the high-resolution achieved by the
HERFD technique and quantum chemical calculations on large
DFT-optimized cluster models of the OEC.
The relative intensities of the pre-edge peaks change clearly
with the S-state progression. Using quantum mechanical
results, we compared different structural variants of each
state on the basis of the difference spectra, the sum of
intensities, and the IWAEs. Our results demonstrate that this
QM-supported analysis of the pre-edge is highly sensitive and
can discriminate even between structures with minor differ-
ences such as terminal water ligands in different protonation
states. The best fitting model for the S1 state is one with
collinear pseudo-Jahn−Teller axes on Mn1(III) and Mn4(III)
ions and with both the W1 and W2 terminal ligands of the
Mn4 ion in the aquo form. The S2 state is most consistent with
an “open-cubane” conformation with an Mn1(III) ion and
both Mn4-bound W1 and W2 in the aquo form. In the S0 state,
protonation of the O4 gives a better fitting with the XAS
Figure 11. Assignment of the calculated pre-edge XAS spectrum based on the NTOs associated with the transitions for models S3
A,W, S3
oxyl−oxo, and
S3
peroxo. In the stick spectra, transitions from 1s core orbitals of Mn1 are shown in green sticks, of Mn2 in red, of Mn3 in blue, and of Mn4 in
orange. NTOs for transitions from 1s core orbitals of Mn1 are shown in green boxes, of Mn3 in blue, and of Mn4 in orange box.
Journal of the American Chemical Society
pubs.acs.org/JACS
Article
https://doi.org/10.1021/jacs.3c06046
J. Am. Chem. Soc. 2023, 145, 25579−25594
25589
